% --- Archon physics formalization source begin ---
% archon:physics
% archon:covers PhyXMiniProblems/problem_phyx_mini_0874.lean
% archon:source-report reports/phyx_mini/problem_phyx_mini_0874.source.json

\chapter{Physics problem phyx\_mini\_0874}
\label{ch:PhyXMiniProblems_problem_phyx_mini_0874}

\paragraph{Problem source.}
\#\# Question

What is the potential difference between A and B?

\#\# Answer choices

- A: A: 12V

- B: B: -14V

- C: C: -200V

- D: D: -70V

\#\# Figure caption (auxiliary; use the image as primary evidence)

The image shows a diagram of an electric field with uniformly spaced, parallel magenta arrows pointing to the right. There are two labeled points, A and B, represented by small black dots.

- **Point A** is positioned near the top-left of the image.
- **Point B** is positioned slightly below and to the right of Point A.

Text in the image includes:
- The electric field strength is labeled as \textbackslash{}( E = 1000 \textbackslash{}, \textbackslash{}text\{V/m\} \textbackslash{}), located to the right of the middle arrow.
- A vertical double-headed arrow on the left side is labeled "3 cm", indicating the vertical distance between the lines.
- A horizontal double-headed arrow at the bottom is labeled "7 cm", indicating the horizontal distance between points A and B.

The parallel arrows suggest the presence of a uniform electric field in the region.

\#\# Dataset metadata

Electrostatics; Physical Model Grounding Reasoning, Spatial Relation Reasoning

\paragraph{Recorded answer/context.}
D: D: -70V

\paragraph{Figure/image path.}
/root/proposal\_for\_physic/science-mango/phyx\_mini\_run/phyx\_data/test\_image/874.png

\paragraph{Formalization target.}
create a compiling Lean file with sorry bodies at `PhyXMiniProblems/problem\_phyx\_mini\_0874.lean`. The Lean declarations must preserve the physical quantities, dimensions or dimensional roles, figure labels, governing-law hypotheses, and final relation expressed by this problem.
Use Mathlib/Physlib names found through LeanExplore where available. If a physics API is missing, introduce faithful local abstractions rather than scalar placeholder aliases.

\begin{theorem}[Physics formalization target]
\label{thm:physics:phyx_mini_0874:target}
\lean{PhyXMiniProblems.ProblemPhyXMini0874.problem_phyx_mini_0874}
\uses{def:physics:phyx-mini-0874:phyxminiproblems-problemphyxmini0874-potentialdifferenceinvolts, def:physics:phyx-mini-0874:phyxminiproblems-problemphyxmini0874-uniformelectricfieldsetup, def:physics:phyx-mini-0874:phyxminiproblems-problemphyxmini0874-matchessupplieduniformfieldfigure, def:physics:phyx-mini-0874:phyxminiproblems-problemphyxmini0874-hasdisplayedpointgeometry, def:physics:phyx-mini-0874:phyxminiproblems-problemphyxmini0874-hasphysicaluniformfielddata, def:physics:phyx-mini-0874:phyxminiproblems-problemphyxmini0874-isstaticuniformelectricfield, def:physics:phyx-mini-0874:phyxminiproblems-problemphyxmini0874-satisfiesuniformfieldpotentiallaw}
The assigned autoformalize agent should translate this physics problem into Lean declarations in the covered file, with theorem and lemma proof bodies written as `by sorry`.
\end{theorem}
\begin{proof}
This is an autoformalization task, not a proof task. Produce statements that can later be proved by the physics prover without weakening the physical model.
\end{proof}

% --- Archon declaration topology begin ---
\section{Declaration topology}
\begin{definition}[Spatial Readout]
  \label{def:physics:phyx-mini-0874:phyxminiproblems-problemphyxmini0874-spatialreadout}
  \lean{PhyXMiniProblems.ProblemPhyXMini0874.SpatialReadout}
  Two-dimensional coordinate vectors at a coherent-SI readout boundary
\end{definition}
\begin{proof}
  This is the declared model definition of Spatial Readout.
\end{proof}
\begin{definition}[electric Potential Dimension]
  \label{def:physics:phyx-mini-0874:phyxminiproblems-problemphyxmini0874-electricpotentialdimension}
  \lean{PhyXMiniProblems.ProblemPhyXMini0874.electricPotentialDimension}
  The physical dimension M L² T⁻² C⁻¹ of electric potential
\end{definition}
\begin{proof}
  This is the declared model definition of electric Potential Dimension.
\end{proof}
\begin{definition}[electric Field Dimension]
  \label{def:physics:phyx-mini-0874:phyxminiproblems-problemphyxmini0874-electricfielddimension}
  \lean{PhyXMiniProblems.ProblemPhyXMini0874.electricFieldDimension}
  The physical dimension M L T⁻² C⁻¹ of electric-field strength
\end{definition}
\begin{proof}
  This is the declared model definition of electric Field Dimension.
\end{proof}
\begin{definition}[Length Quantity]
  \label{def:physics:phyx-mini-0874:phyxminiproblems-problemphyxmini0874-lengthquantity}
  \lean{PhyXMiniProblems.ProblemPhyXMini0874.LengthQuantity}
  A nonnegative physical length
\end{definition}
\begin{proof}
  This is the declared model definition of Length Quantity.
\end{proof}
\begin{definition}[Position Quantity]
  \label{def:physics:phyx-mini-0874:phyxminiproblems-problemphyxmini0874-positionquantity}
  \lean{PhyXMiniProblems.ProblemPhyXMini0874.PositionQuantity}
  \uses{def:physics:phyx-mini-0874:phyxminiproblems-problemphyxmini0874-spatialreadout}
  A signed two-dimensional physical position
\end{definition}
\begin{proof}
  This is the declared model definition of Position Quantity.
\end{proof}
\begin{definition}[Electric Field Vector Quantity]
  \label{def:physics:phyx-mini-0874:phyxminiproblems-problemphyxmini0874-electricfieldvectorquantity}
  \lean{PhyXMiniProblems.ProblemPhyXMini0874.ElectricFieldVectorQuantity}
  \uses{def:physics:phyx-mini-0874:phyxminiproblems-problemphyxmini0874-spatialreadout, def:physics:phyx-mini-0874:phyxminiproblems-problemphyxmini0874-electricfielddimension}
  A signed two-dimensional physical electric-field vector
\end{definition}
\begin{proof}
  This is the declared model definition of Electric Field Vector Quantity.
\end{proof}
\begin{definition}[Electric Potential Quantity]
  \label{def:physics:phyx-mini-0874:phyxminiproblems-problemphyxmini0874-electricpotentialquantity}
  \lean{PhyXMiniProblems.ProblemPhyXMini0874.ElectricPotentialQuantity}
  \uses{def:physics:phyx-mini-0874:phyxminiproblems-problemphyxmini0874-electricpotentialdimension}
  A signed physical electric potential
\end{definition}
\begin{proof}
  This is the declared model definition of Electric Potential Quantity.
\end{proof}
\begin{definition}[Potential Difference Quantity]
  \label{def:physics:phyx-mini-0874:phyxminiproblems-problemphyxmini0874-potentialdifferencequantity}
  \lean{PhyXMiniProblems.ProblemPhyXMini0874.PotentialDifferenceQuantity}
  \uses{def:physics:phyx-mini-0874:phyxminiproblems-problemphyxmini0874-electricpotentialdimension}
  A signed, oriented physical electric-potential difference
\end{definition}
\begin{proof}
  This is the declared model definition of Potential Difference Quantity.
\end{proof}
\begin{definition}[length Readout]
  \label{def:physics:phyx-mini-0874:phyxminiproblems-problemphyxmini0874-lengthreadout}
  \lean{PhyXMiniProblems.ProblemPhyXMini0874.lengthReadout}
  \uses{def:physics:phyx-mini-0874:phyxminiproblems-problemphyxmini0874-lengthquantity}
  Read a physical length in a selected PhysLean length unit
\end{definition}
\begin{proof}
  This is the declared model definition of length Readout.
\end{proof}
\begin{definition}[length In Meters]
  \label{def:physics:phyx-mini-0874:phyxminiproblems-problemphyxmini0874-lengthinmeters}
  \lean{PhyXMiniProblems.ProblemPhyXMini0874.lengthInMeters}
  \uses{def:physics:phyx-mini-0874:phyxminiproblems-problemphyxmini0874-lengthquantity, def:physics:phyx-mini-0874:phyxminiproblems-problemphyxmini0874-lengthreadout}
  Read a physical length in coherent-SI metres
\end{definition}
\begin{proof}
  This is the declared model definition of length In Meters.
\end{proof}
\begin{definition}[length In Centimeters]
  \label{def:physics:phyx-mini-0874:phyxminiproblems-problemphyxmini0874-lengthincentimeters}
  \lean{PhyXMiniProblems.ProblemPhyXMini0874.lengthInCentimeters}
  \uses{def:physics:phyx-mini-0874:phyxminiproblems-problemphyxmini0874-lengthquantity, def:physics:phyx-mini-0874:phyxminiproblems-problemphyxmini0874-lengthreadout}
  Read a physical length in centimetres
\end{definition}
\begin{proof}
  This is the declared model definition of length In Centimeters.
\end{proof}
\begin{definition}[position In Meters]
  \label{def:physics:phyx-mini-0874:phyxminiproblems-problemphyxmini0874-positioninmeters}
  \lean{PhyXMiniProblems.ProblemPhyXMini0874.positionInMeters}
  \uses{def:physics:phyx-mini-0874:phyxminiproblems-problemphyxmini0874-spatialreadout, def:physics:phyx-mini-0874:phyxminiproblems-problemphyxmini0874-positionquantity}
  Read a physical position in coherent-SI metres
\end{definition}
\begin{proof}
  This is the declared model definition of position In Meters.
\end{proof}
\begin{definition}[electric Field In Volts Per Meter]
  \label{def:physics:phyx-mini-0874:phyxminiproblems-problemphyxmini0874-electricfieldinvoltspermeter}
  \lean{PhyXMiniProblems.ProblemPhyXMini0874.electricFieldInVoltsPerMeter}
  \uses{def:physics:phyx-mini-0874:phyxminiproblems-problemphyxmini0874-spatialreadout, def:physics:phyx-mini-0874:phyxminiproblems-problemphyxmini0874-electricfieldvectorquantity}
  Read an electric-field vector in coherent-SI volts per metre
\end{definition}
\begin{proof}
  This is the declared model definition of electric Field In Volts Per Meter.
\end{proof}
\begin{definition}[electric Potential In Volts]
  \label{def:physics:phyx-mini-0874:phyxminiproblems-problemphyxmini0874-electricpotentialinvolts}
  \lean{PhyXMiniProblems.ProblemPhyXMini0874.electricPotentialInVolts}
  \uses{def:physics:phyx-mini-0874:phyxminiproblems-problemphyxmini0874-electricpotentialquantity}
  Read an electric potential in coherent-SI volts
\end{definition}
\begin{proof}
  This is the declared model definition of electric Potential In Volts.
\end{proof}
\begin{definition}[potential Difference In Volts]
  \label{def:physics:phyx-mini-0874:phyxminiproblems-problemphyxmini0874-potentialdifferenceinvolts}
  \lean{PhyXMiniProblems.ProblemPhyXMini0874.potentialDifferenceInVolts}
  \uses{def:physics:phyx-mini-0874:phyxminiproblems-problemphyxmini0874-potentialdifferencequantity}
  Read a signed, oriented electric-potential difference in volts
\end{definition}
\begin{proof}
  This is the declared model definition of potential Difference In Volts.
\end{proof}
\begin{lemma}[length In Centimeters eq one Hundred mul length In Meters]
  \label{lem:physics:phyx-mini-0874:phyxminiproblems-problemphyxmini0874-lengthincentimeters-eq-onehundred-mul-lengthinmeters}
  \lean{PhyXMiniProblems.ProblemPhyXMini0874.lengthInCentimeters_eq_oneHundred_mul_lengthInMeters}
  \uses{def:physics:phyx-mini-0874:phyxminiproblems-problemphyxmini0874-lengthquantity, def:physics:phyx-mini-0874:phyxminiproblems-problemphyxmini0874-lengthinmeters, def:physics:phyx-mini-0874:phyxminiproblems-problemphyxmini0874-lengthincentimeters}
  The centimetre-to-metre conversion used by the displayed geometry
\end{lemma}
\begin{proof}
  The conclusion follows from the governing relations and figure readouts named in the statement of length In Centimeters eq one Hundred mul length In Meters.
\end{proof}
\begin{definition}[Diagram Point]
  \label{def:physics:phyx-mini-0874:phyxminiproblems-problemphyxmini0874-diagrampoint}
  \lean{PhyXMiniProblems.ProblemPhyXMini0874.DiagramPoint}
  The two labelled observation points in the supplied image
\end{definition}
\begin{proof}
  This is the declared model definition of Diagram Point.
\end{proof}
\begin{definition}[Field Arrow Label]
  \label{def:physics:phyx-mini-0874:phyxminiproblems-problemphyxmini0874-fieldarrowlabel}
  \lean{PhyXMiniProblems.ProblemPhyXMini0874.FieldArrowLabel}
  The six distinct magenta field-arrow glyphs visible in the image
\end{definition}
\begin{proof}
  This is the declared model definition of Field Arrow Label.
\end{proof}
\begin{definition}[Diagram Direction]
  \label{def:physics:phyx-mini-0874:phyxminiproblems-problemphyxmini0874-diagramdirection}
  \lean{PhyXMiniProblems.ProblemPhyXMini0874.DiagramDirection}
  Qualitative arrow directions in the diagram plane
\end{definition}
\begin{proof}
  This is the declared model definition of Diagram Direction.
\end{proof}
\begin{definition}[Printed Electric Field Unit]
  \label{def:physics:phyx-mini-0874:phyxminiproblems-problemphyxmini0874-printedelectricfieldunit}
  \lean{PhyXMiniProblems.ProblemPhyXMini0874.PrintedElectricFieldUnit}
  Unit glyphs available for the printed electric-field label
\end{definition}
\begin{proof}
  This is the declared model definition of Printed Electric Field Unit.
\end{proof}
\begin{definition}[Printed Length Unit]
  \label{def:physics:phyx-mini-0874:phyxminiproblems-problemphyxmini0874-printedlengthunit}
  \lean{PhyXMiniProblems.ProblemPhyXMini0874.PrintedLengthUnit}
  Unit glyphs available for the printed distance labels
\end{definition}
\begin{proof}
  This is the declared model definition of Printed Length Unit.
\end{proof}
\begin{definition}[Uniform Electric Field Figure]
  \label{def:physics:phyx-mini-0874:phyxminiproblems-problemphyxmini0874-uniformelectricfieldfigure}
  \lean{PhyXMiniProblems.ProblemPhyXMini0874.UniformElectricFieldFigure}
  \uses{def:physics:phyx-mini-0874:phyxminiproblems-problemphyxmini0874-diagrampoint, def:physics:phyx-mini-0874:phyxminiproblems-problemphyxmini0874-fieldarrowlabel, def:physics:phyx-mini-0874:phyxminiproblems-problemphyxmini0874-diagramdirection, def:physics:phyx-mini-0874:phyxminiproblems-problemphyxmini0874-printedelectricfieldunit, def:physics:phyx-mini-0874:phyxminiproblems-problemphyxmini0874-printedlengthunit}
  Literal presentation data transcribed from 874.png. Its scalar fields are unit-labelled raster readouts, not replacements for physical quantities
\end{definition}
\begin{proof}
  This is the declared model definition of Uniform Electric Field Figure.
\end{proof}
\begin{definition}[Uniform Electric Field Setup]
  \label{def:physics:phyx-mini-0874:phyxminiproblems-problemphyxmini0874-uniformelectricfieldsetup}
  \lean{PhyXMiniProblems.ProblemPhyXMini0874.UniformElectricFieldSetup}
  \uses{def:physics:phyx-mini-0874:phyxminiproblems-problemphyxmini0874-spatialreadout, def:physics:phyx-mini-0874:phyxminiproblems-problemphyxmini0874-lengthquantity, def:physics:phyx-mini-0874:phyxminiproblems-problemphyxmini0874-positionquantity, def:physics:phyx-mini-0874:phyxminiproblems-problemphyxmini0874-electricfieldvectorquantity, def:physics:phyx-mini-0874:phyxminiproblems-problemphyxmini0874-electricpotentialquantity, def:physics:phyx-mini-0874:phyxminiproblems-problemphyxmini0874-potentialdifferencequantity, def:physics:phyx-mini-0874:phyxminiproblems-problemphyxmini0874-diagrampoint, def:physics:phyx-mini-0874:phyxminiproblems-problemphyxmini0874-uniformelectricfieldfigure}
  Independent physical quantities for the uniform-field configuration. potentialDifference start finish is oriented as V(finish) - V(start)
\end{definition}
\begin{proof}
  This is the declared model definition of Uniform Electric Field Setup.
\end{proof}
\begin{definition}[Matches Supplied Uniform Field Figure]
  \label{def:physics:phyx-mini-0874:phyxminiproblems-problemphyxmini0874-matchessupplieduniformfieldfigure}
  \lean{PhyXMiniProblems.ProblemPhyXMini0874.MatchesSuppliedUniformFieldFigure}
  \uses{def:physics:phyx-mini-0874:phyxminiproblems-problemphyxmini0874-lengthincentimeters, def:physics:phyx-mini-0874:phyxminiproblems-problemphyxmini0874-electricfieldinvoltspermeter, def:physics:phyx-mini-0874:phyxminiproblems-problemphyxmini0874-uniformelectricfieldsetup}
  Primary-image evidence and its calibration to independent physical quantities. These assumptions contain the printed 1000, 7, and 3, but do not contain the requested potential difference
\end{definition}
\begin{proof}
  This is the declared model definition of Matches Supplied Uniform Field Figure.
\end{proof}
\begin{definition}[Has Displayed Point Geometry]
  \label{def:physics:phyx-mini-0874:phyxminiproblems-problemphyxmini0874-hasdisplayedpointgeometry}
  \lean{PhyXMiniProblems.ProblemPhyXMini0874.HasDisplayedPointGeometry}
  \uses{def:physics:phyx-mini-0874:phyxminiproblems-problemphyxmini0874-lengthinmeters, def:physics:phyx-mini-0874:phyxminiproblems-problemphyxmini0874-positioninmeters, def:physics:phyx-mini-0874:phyxminiproblems-problemphyxmini0874-uniformelectricfieldsetup}
  The diagram directions form an orthonormal frame, and B is the displayed horizontal distance right and vertical distance below A. This is geometry from the image rather than a potential conclusion
\end{definition}
\begin{proof}
  This is the declared model definition of Has Displayed Point Geometry.
\end{proof}
\begin{definition}[Has Physical Uniform Field Data]
  \label{def:physics:phyx-mini-0874:phyxminiproblems-problemphyxmini0874-hasphysicaluniformfielddata}
  \lean{PhyXMiniProblems.ProblemPhyXMini0874.HasPhysicalUniformFieldData}
  \uses{def:physics:phyx-mini-0874:phyxminiproblems-problemphyxmini0874-lengthinmeters, def:physics:phyx-mini-0874:phyxminiproblems-problemphyxmini0874-uniformelectricfieldsetup}
  Positivity conditions for the nondegenerate distances and field label
\end{definition}
\begin{proof}
  This is the declared model definition of Has Physical Uniform Field Data.
\end{proof}
\begin{definition}[Is Static Uniform Electric Field]
  \label{def:physics:phyx-mini-0874:phyxminiproblems-problemphyxmini0874-isstaticuniformelectricfield}
  \lean{PhyXMiniProblems.ProblemPhyXMini0874.IsStaticUniformElectricField}
  \uses{def:physics:phyx-mini-0874:phyxminiproblems-problemphyxmini0874-electricfieldinvoltspermeter, def:physics:phyx-mini-0874:phyxminiproblems-problemphyxmini0874-uniformelectricfieldsetup}
  The PhysLean spacetime field is static and uniform, and agrees everywhere with the coherent-SI readout of the independent dimensionful field vector
\end{definition}
\begin{proof}
  This is the declared model definition of Is Static Uniform Electric Field.
\end{proof}
\begin{definition}[Satisfies Uniform Field Potential Law]
  \label{def:physics:phyx-mini-0874:phyxminiproblems-problemphyxmini0874-satisfiesuniformfieldpotentiallaw}
  \lean{PhyXMiniProblems.ProblemPhyXMini0874.SatisfiesUniformFieldPotentialLaw}
  \uses{def:physics:phyx-mini-0874:phyxminiproblems-problemphyxmini0874-positioninmeters, def:physics:phyx-mini-0874:phyxminiproblems-problemphyxmini0874-electricfieldinvoltspermeter, def:physics:phyx-mini-0874:phyxminiproblems-problemphyxmini0874-electricpotentialinvolts, def:physics:phyx-mini-0874:phyxminiproblems-problemphyxmini0874-potentialdifferenceinvolts, def:physics:phyx-mini-0874:phyxminiproblems-problemphyxmini0874-uniformelectricfieldsetup}
  General electrostatic potential laws for every ordered pair of labelled points. The first law fixes the orientation of potential difference. The second is the constant-field relation V(finish) - V(start) = -E · (r(finish) - r(start))
\end{definition}
\begin{proof}
  This is the declared model definition of Satisfies Uniform Field Potential Law.
\end{proof}
\begin{definition}[Answer Choice]
  \label{def:physics:phyx-mini-0874:phyxminiproblems-problemphyxmini0874-answerchoice}
  \lean{PhyXMiniProblems.ProblemPhyXMini0874.AnswerChoice}
  The four answer labels supplied with the problem
\end{definition}
\begin{proof}
  This is the declared model definition of Answer Choice.
\end{proof}
\begin{definition}[displayed Potential Difference In Volts]
  \label{def:physics:phyx-mini-0874:phyxminiproblems-problemphyxmini0874-answerchoice-displayedpotentialdifferenceinvolts}
  \lean{PhyXMiniProblems.ProblemPhyXMini0874.AnswerChoice.displayedPotentialDifferenceInVolts}
  \uses{def:physics:phyx-mini-0874:phyxminiproblems-problemphyxmini0874-answerchoice}
  The signed potential difference printed for each answer choice, in volts
\end{definition}
\begin{proof}
  This is the declared model definition of displayed Potential Difference In Volts.
\end{proof}
\begin{definition}[recorded Dataset Answer]
  \label{def:physics:phyx-mini-0874:phyxminiproblems-problemphyxmini0874-recordeddatasetanswer}
  \lean{PhyXMiniProblems.ProblemPhyXMini0874.recordedDatasetAnswer}
  \uses{def:physics:phyx-mini-0874:phyxminiproblems-problemphyxmini0874-answerchoice}
  The answer label recorded in the source dataset
\end{definition}
\begin{proof}
  This is the declared model definition of recorded Dataset Answer.
\end{proof}
% --- Archon declaration topology end ---
% --- Archon physics formalization source end ---
